% Appendix: what failed for each model and what fixed it. Numbers from RESULTS.md (one seed unless stated).
\section{What Failed and How It Was Fixed}
\label{sec:failures}

This appendix records, model by model, every failure we hit, the evidence for it, the diagnosis, and what
fixed it. Unresolved failures are reported as such. All runs fine-tune the released TinyStories NCM
(88.6M parameters, 11 target levels, 641 codes per 128-token chunk) on $M{=}4$ chunks (512 tokens) for 1{,}500
steps of 8 documents. NLL is measured on all 1{,}228 validation documents with padding masked. genPPL uses
GPT-2 Large, $2$ seeds $\times$ 256 samples, temperature 1.0. The \emph{boundary ratio} is the genPPL of 64-token
windows across a chunk seam divided by that of windows inside a chunk (real text: 1.01).

\subsection{Shared pipeline (all models)}

\paragraph{F1. Fine-tuning destroyed free generation.}
\emph{Evidence.} The first $M{=}4$ runs (clean inputs, lr $10^{-4}$) reached good teacher-forced NLL, but
genPPL was 47.8--53.5 for every model. Fine-tuned B0 also scored 47.7 on \emph{single} 128-token chunks, against
16.5 for the released model, with no chunking involved. On story openings, fine-tuned B0 beat the released
model on NLL (1.179 vs 1.558, better on 100\% of documents) yet generated 3$\times$ worse.
\emph{Diagnosis.} The released generator is trained with \emph{conditioning corruption}: a fraction of input
codes is replaced by random codes from the same codebook, at a rate decreasing linearly over levels (mean 0.5;
NCM paper, Sec.~3.4 and App.~F.3). Clean fine-tuning makes the model trust its conditioning, so at sampling time
its own errors compound. We ruled out the alternatives:
\begin{itemize}[nosep]
  \item train and generation inputs are bit-identical at every level;
  \item the tokenizer is unchanged by fine-tuning;
  \item sampling settings match the released evaluation;
  \item a train-save-load-generate control at lr $10^{-12}$ gives 16.6.
\end{itemize}
\emph{Fix.} Uniform code replacement with a linear 0.75$\to$0.25 schedule (targets clean) and lr
$3{\times}10^{-5}$, chosen by a single-chunk gate (Table~\ref{tab:gate}).
\emph{Outcome.} Single-chunk genPPL 47.7$\to$20.7. A residual gap of about 4--5 genPPL remained (20.5 plateau vs
16.5). It did not respond to learning rate ($10^{-5}$ to $3{\times}10^{-5}$), to a 4$\times$ larger batch, or to
per-level loss weighting (20.9).
\emph{Diagnosis of the residual.} The single-chunk test trained on the openings of stories longer than 384 tokens
(our $M{=}4$ set, 4.7\% of TinyStories). These are a different style of story: none opens with ``Once upon a
time'', against 58\% of all stories and 70\% of the released model's samples. Trained on 60k unfiltered openings
instead, the same recipe scores 17.0 (entropy 3.98, as released); adding a KL term that anchors the model to the
frozen released model gives 16.1, the released level within noise. The anchor alone, on the long-story openings,
gives 19.5. \textbf{Mostly resolved:} the residual was largely the data subset, not damage from fine-tuning. The
$M{=}4$ models must train on long stories, so their absolute genPPL carries this style offset (real long-story
openings score 1.10$\times$ higher than all openings); the cost is the same for every variant, so comparisons
between variants remain fair.

\begin{table}[h]
\centering
\caption{Single-chunk gate: 300 steps on story openings, genPPL of 256 samples (released model: 16.5).}
\label{tab:gate}
\begin{tabular}{lcc}
\toprule
recipe & genPPL & entropy \\
\midrule
lr $10^{-4}$, clean & 41.0 & 3.99 \\
lr $10^{-5}$, clean & 34.1 & 3.98 \\
lr $10^{-4}$, corruption mean 0.25 & 24.4 & 3.94 \\
lr $10^{-4}$, corruption mean 0.5 (linear) & 22.2 & 3.92 \\
lr $10^{-4}$, corruption 0.5 (flat) & 21.8 & 3.92 \\
\textbf{lr $3{\times}10^{-5}$, corruption mean 0.5 (chosen)} & \textbf{20.7} & 3.92 \\
\bottomrule
\end{tabular}
\end{table}

\paragraph{F2. Padding inflated the agreement between models.}
37\% of chunk-4 codes are padding: stories end before token 512. Padding is equally easy for every model and
diluted the differences. \emph{Fix:} per-document lengths are saved at encoding time and padding codes are
masked in every NLL. A per-document re-encoding check confirms the mask lines up with the stored codes.
\emph{Outcome:} the B0$\to$B2 gap grew from 0.015 to 0.020 nats in the clean-recipe runs; all reported NLL is
padding-free.

\paragraph{F3. Memory: the training graph did not fit in 8\,GB.}
The single dense-mask training sequence ($M{=}4$, all levels) exceeds 8\,GB. \emph{Fix:} a per-level training
path (one short sequence per level, built from the generation layouts). Its loss equals the dense path's to within
$10^{-6}$ for $G\in\{0,1,3,6,11\}$, and gradient checkpointing is used throughout.

\paragraph{F4. Mixed precision.}
On the T4 (no bf16), the quantizer's convolutions produced NaN in fp16, and training continued silently on NaN
loss. \emph{Fix:} the quantizer always runs in fp32; training aborts on any non-finite loss; the pipeline greps
every stage log for NaN. Separately, the bf16 quantizer loses precision even on the RTX~5050 (loss 3.116 vs
2.965 in fp32). This may explain why our reproduction of the released model scores 16.38 against the paper's
15.41; \textbf{a test is pending}.

\paragraph{F5. Evaluation running out of memory.}
Generation at $M{=}6$ and plan search (8 plans) ran out of memory at batch 64. \emph{Fix:} every evaluation
retries at half the batch size. Batch size changes only speed, not the samples' distribution.

\subsection{B0: independent chunks ($G{=}0$)}
\emph{Failure (by design):} no information crosses chunks. With the fixed recipe it has the worst seams
(boundary ratio 1.52) and frequent character-name drift between chunks; genPPL is 24.1. \emph{No fix:} B0 is the
lower bound that the plan must beat.

\subsection{Ours: Block-NCM ($G{=}3$, plan over levels 1--3)}
\paragraph{O1. The plan alone did not fix the seams.}
\emph{Evidence.} Ours has the best NLL of the three baselines (1.7438 vs 1.7486 B0 and 1.7780 B2; $-0.0049$ nats
vs B0, 95\% CI $[-0.0056,-0.0042]$). But it closed only about 40\% of the B0$\to$B2 gap in generation: genPPL
22.0 vs 19.0, boundary ratio 1.39 vs 1.17.
\emph{Diagnosis.} A level-3 plan carries topic and characters. Seams need token-level continuity, though, and
that is decided by the local fine levels, which cannot see the neighbouring chunk.
\emph{Fix: halo attention.} Each local level also attends to the inputs at the last $h$ positions of the previous
chunk and the first $h$ of the next. These inputs come from coarser levels that are already generated, so this
adds no sequential passes. It costs $2h$ extra positions per chunk per local level (about 10\% compute for $h{=}4$).
\emph{Outcome} ($h{=}4$):
\begin{itemize}[nosep]
  \item boundary ratio 1.39$\to$1.19, closing 95\% of the B0$\to$B2 gap;
  \item best NLL of all runs (1.7361);
  \item genPPL 22.0$\to$21.0.
\end{itemize}
Plan search (8 sampled plans, keep the most likely) adds a further gain: 20.0. Wider halos help a little:
$h=2,4,8$ give genPPL 21.4, 21.0 and 20.4, with boundary ratios 1.22, 1.19 and 1.17.

\paragraph{O2. Recomputing the plan made training slow.}
The reference implementation recomputes the global plan for every local level.
\emph{Fix:} compute the plan's keys and values once and cache them. This is exact (the losses match), and a
training step is 2.0$\times$ faster on a T4 (1.21\,s vs 2.43\,s).

\paragraph{O3. Generation quality is still below the released model.}
Samples keep a coherent plot across all four chunks, but grammar slips are frequent. That is a consequence of F1
and of model size. Plan-only adapters did not help (below). Under test: lower sampling temperature, and recipe
changes (unfiltered openings, per-level loss weighting, a KL anchor to the released model).
\textbf{Open.}

\subsection{Ours with more global levels ($G{=}5$, $G{=}6$)}
More global levels improve generation at the training length: genPPL 20.0 ($G{=}5$) and 19.6 ($G{=}6$) vs 19.0
for B2. But they \emph{fail to extrapolate}. At $M{=}6$ (trained on $M{=}4$), $G{=}6$ degrades to 30.7 and
$G{=}5$ to 23.8, while $G{=}3$ holds (21.8). $G{=}6$ also loses on the fine levels (level 6 NLL 1.86 vs 1.68).
\emph{Fix: add halo.} With halo $h{=}4$, $G{=}5$ goes from 19.05 at $M{=}4$ to 20.5 at $M{=}6$ (means of two seeds),
the best $M{=}6$ result of any model; B2 goes from 18.7 to 32.3. At $M{=}4$, B2 leads by about 0.4 genPPL and the
seams tie (boundary ratio 1.14 vs 1.13). $G{=}6$ with halo generates best at $M{=}4$ (18.7) but still degrades at $M{=}6$ (23.6).
\emph{Resolution:} $G{=}5$ with halo is the final model; $G\ge 6$ learns the training length even with halo.

\subsection{B2: fully global ($G{=}11$)}
B2 generates best at $M{=}4$ but has the \emph{worst} NLL of all runs (1.7780), losing on levels 6--11 and on
chunk 4. More importantly, it \textbf{collapses beyond the training length}: genPPL 19.0$\to$33.5 at $M{=}6$, and
18.4$\to$31.0 in a second training seed, with high sampling variance (sd 2.3 and 1.2). Global attention across all chunks appears to fit the training length. This is the
baseline's failure, not ours, and it motivates Block-NCM: B2's advantage exists only at the length it was
trained on.

\subsection{Plan-only adapters (frozen generator)}
Goal: keep the released model's fluency (F1) by freezing the generator and training only LoRA adapters on the
global passes, plus a learned bias on attention to the plan.
\begin{itemize}[nosep]
  \item \emph{Crash:} a float attention mask with gradients fails in the fused GPU attention kernel.
        \emph{Fix:} implement the plan bias as an extra key dimension.
  \item \emph{Bias never opened:} Adam moves a parameter by about one learning-rate step per update, so a bias
        initialised at $-8$ with lr $10^{-4}$ would need $10^5$ steps. \emph{Fix:} initialise at $-6$, with its
        own lr $10^{-2}$.
  \item \emph{Leak:} adapters also changed the prefix and the chunk embeddings of local levels, so the model was no
        longer ``plan-only''. \emph{Fix:} adapters skip the prefix, and chunk embeddings apply to global levels
        only. Seven unit checks verify this.
\end{itemize}
\emph{Outcome:} \textbf{failed}. The adapters have the worst NLL of all fine-tuned models (1.834 vs 1.744 for Ours;
chunk 4: 2.22 vs 1.79), because the frozen local levels never learn chunks 2--4. They also generate worst
(genPPL 28.5 vs 22.0 at $M{=}4$), so freezing did not preserve the released model's fluency. Full fine-tuning stays
the method.

\subsection{Layer-drop refiner (negative result, moved to future work)}
Goal: cut serving cost by running local levels with 6 of the 12 layers, with each kept layer reading the plan
cache of the matching global layer.
\begin{itemize}[nosep]
  \item \emph{Untrained:} the best 6-layer pattern ($[0,1,2,9,10,11]$) loses 4.0 nats on levels 4--11.
  \item \emph{Fix attempted:} distillation from the full-depth model (KL weight 1.0, 1{,}500 steps). This
        recovered most of the loss, but not enough:
        \begin{itemize}[nosep]
          \item +0.50 nats with genPPL 21.6$\to$41.6 (levels 4--11, 1.55$\times$ faster per sample);
          \item +0.30 nats with genPPL 32.0 (levels 7--11, 1.38$\times$ faster).
        \end{itemize}
\end{itemize}
\textbf{Not fixed.} Depth reduction is left to follow-up work. Block-NCM's serving-cost claims are therefore
restricted to low batch sizes (Sec.~\ref{sec:complexity}).

\subsection{Summary}
\begin{table}[h]
\centering
\small
\begin{tabular}{p{2.6cm}p{4.2cm}p{4.2cm}p{2.2cm}}
\toprule
model / part & failure & fix & status \\
\midrule
all (recipe) & clean fine-tuning: genPPL 47.7 on single chunks & corruption 0.5, lr $3{\times}10^{-5}$; residual: unfiltered data + KL anchor & 16.1 (released 16.5), mostly fixed \\
all (metrics) & padding diluted NLL & mask padding per document & fixed \\
all (memory) & dense graph $>$ 8\,GB & per-level path ($<10^{-6}$ equal) & fixed \\
all (fp16) & quantizer NaN, silent & fp32 quantizer, abort on NaN & fixed \\
Ours $G{=}3$ & seams 1.39 & halo attention, $h{=}4$ & 1.19, fixed \\
Ours $G{=}3$ & plan recomputed per level & KV-cached plan (exact) & 2.0$\times$, fixed \\
Ours $G{=}5,6$ & degrade at $M{=}6$ (23.8, 30.7) & halo; $G{=}5$ & $G{=}5$: 20.5 (2 seeds), fixed; $G{=}6$: 23.6 \\
B2 & collapse at $M{=}6$ (33.5, 31.0) & none (baseline) & finding \\
plan-only adapters & mask crash; bias stuck; leak; then NLL 1.834, genPPL 28.5 & key-dim bias; own lr; skip prefix & \textbf{failed} \\
layer-drop refiner & +0.3--0.5 nats, genPPL $\times$1.5--2 & distillation (insufficient) & \textbf{failed} \\
\bottomrule
\end{tabular}
\end{table}
